\documentclass[10pt]{article}
\usepackage[T1]{fontenc}
\usepackage[utf8]{inputenc}
\usepackage[french]{babel}
\usepackage{amsmath,amssymb,mathrsfs,mathpazo}
\usepackage{geometry}
\geometry{a4paper,margin=2cm}
\usepackage{graphicx,array,tabularx,tabularray,colortbl,multirow,multicol,booktabs,enumitem,pifont,ulem,fancybox,diagbox,cancel}
\usepackage{pstricks,pstricks-add,pst-plot,pst-tree,pst-node,pst-text,pst-func,pst-math,pst-bspline,pst-3dplot,pst-eucl,pst-solides3d}
\let\pscalculate\undefined
\usepackage{pst-bezier}
\usepackage{tikz,pgfplots,tkz-tab}
\usetikzlibrary{arrows,arrows.meta,patterns,calc,positioning,trees,decorations.markings,shapes.geometric,shapes.misc,shapes.arrows,intersections,decorations.pathmorphing,decorations.pathreplacing,backgrounds,shadows,babel}
\pgfplotsset{compat=1.18}
\usepgfplotslibrary{fillbetween}
\usepackage[np]{numprint}
\usepackage{listings,makecell,pas-tableur,twoopt}
\let\barre\undefined
\usepackage[pstricks]{tablvar}
\definecolor{darkgreen}{RGB}{0,100,0}
\providecommand{\red}{\color{red}}
\providecommand{\blue}{\color{blue}}
\providecommand{\green}{\color{green}}
\providecommand{\black}{\color{black}}
\providecommand{\white}{\color{white}}
\providecommand{\gray}{\color{gray}}
\providecommand{\cyan}{\color{cyan}}
\providecommand{\magenta}{\color{magenta}}
\providecommand{\yellow}{\color{yellow}}
\providecommand{\violet}{\color{violet}}
\providecommand{\orange}{\color{orange}}
\providecommand{\euro}{\texteuro}
\providecommand{\EUR}{\texteuro}
\providecommand{\ieme}{\textsuperscript{e}}
\providecommand{\degres}{\ensuremath{^\circ}}
\providecommand{\degree}{\ensuremath{^\circ}}
\providecommand{\vv}[1]{\overrightarrow{#1}}
\providecommand{\Coord}[1]{#1}
\providecommand{\decosix}{}
\providecommand{\starredbullet}{\ensuremath{\bullet}}
\pagestyle{empty}
\setlength{\parindent}{0pt}
\usepackage{adjustbox}\DeclareUnicodeCharacter{2236}{:}\DeclareUnicodeCharacter{2190}{\ensuremath{\leftarrow}}
\providecommand{\R}{\mathbb{R}}
\providecommand{\N}{\mathbb{N}}
\providecommand{\Z}{\mathbb{Z}}
\providecommand{\Q}{\mathbb{Q}}
\providecommand{\D}{\mathbb{D}}
\providecommand{\C}{\mathbb{C}}
\providecommand{\e}{{\mathrm{e}}}
\providecommand{\vect}[1]{\overrightarrow{#1}}
\providecommand{\vectt}[1]{\overrightarrow{\mathrm{#1}}}
\providecommand{\vvt}[1]{\overrightarrow{\mathrm{#1}}}
\providecommand{\np}[2][]{{#2~#1}}
\providecommand{\vv}[1]{{\overrightarrow{#1}}}
\renewcommand{\d}{{\mathrm{d}}}
\providecommand{\degres}{{^\circ}}
\providecommand{\degre}{{^\circ}}
\providecommand{\blue}{\color{blue}}
\providecommand{\red}{\color{red}}
\providecommand{\enskip}{\quad}
\providecommand{\llbracket}{\lbrack\!\lbrack}
\providecommand{\rrbracket}{\rbrack\!\rbrack}
\providecommand{\euro}{{\text{€}}}
\providecommand{\no}{{\text{n°}}}
\providecommand{\interval}[2]{{[#1;#2]}}
\providecommand{\intervFO}[2]{{[#1;#2[}}
\providecommand{\dotfill}{\ldots}

\providecommand{\ds}{\displaystyle}
\providecommand{\textding}[1]{#1}
\providecommand{\Oij}{\(\mathrm{O};\vec{\imath},\vec{\jmath}\)}
\providecommand{\Oijk}{\(\mathrm{O};\vec{\imath},\vec{\jmath},\vec{k}\)}
\providecommand{\Ouv}{\(\mathrm{O};\vec{u},\vec{v}\)}
\providecommand{\intervFF}[2]{\ensuremath{[#1;#2]}}
\renewcommand{\euro}{\ifmmode\text{\texteuro}\else\texteuro\fi}
\DeclareUnicodeCharacter{00B0}{\ifmmode^{\circ}\else\textdegree\fi}
\DeclareUnicodeCharacter{20AC}{\ifmmode\text{\texteuro}\else\texteuro\fi}
% Énoncés : Éducation nationale. Transcription et corrigés : APMEP, Denis Vergès et contributeurs. Découpage et classement : CoopMaths, Sébastien Lozano, Éric Elter et contributeurs. Adaptation de la présentation : Maths Explorer.
% CC BY-SA : https://coopmaths.fr/www/about/
% https://coopmaths.fr/alea/static/bac/2022/tex/bac_2022_10_sujet2_caledonie_1_cor.tex
\begin{document}
\noindent\textit{APMEP / CoopMaths. CC BY-SA. Présentation adaptée par Maths Explorer.}\par\medskip
\textbf{Principaux domaines abordés :} Probabilités

\medskip























\begin{enumerate}
\item 







	\begin{enumerate}
		\item ~
\begin{center}
\pstree[treemode=R,nodesepA=0pt,nodesepB=2.5pt,treesep = 1cm,levelsep=2.5cm]{\TR{}}
{\pstree{\TR{$D$~}\taput{0,25}}
{\TR{$R$}\taput{0,6}
\TR{$\overline{R}$}\tbput{0,4}
}
\pstree{\TR{$\overline{D}$~}\tbput{0,75}}
{\TR{$R$}\taput{0,35}
\TR{$\overline{R}$}\tbput{0,65}
}}
\end{center}
		\item 
$p\left(\overline{D} \cap R\right) = p\left(\overline{D}\right) \times p_{\overline{D}}\left(R \right) = 0,75 \times 0,35 = \np{0,2625}$.

		\item 
On a de même $p(D \cap R) = p(D) \times p_D(R) = 0,25 \times 0,6 = 0,15$.

D'après la loi des probabilités totales : 

$p(R) = p(D \cap R) + p\left(\overline{D} \cap R\right) =   0,15 + \np{0,2625} = \np{0,4125}$.
		\item 
		Il faut trouver $p_R\left(\overline{D}\right) = \dfrac{p\left(R \cap \overline{D}\right)}{p(R)} = \dfrac{\np{0,2625}}{\np{0,4125}}\approx \np{0,6364}$, soit 0,64 au centième près.
	\end{enumerate}	
\item 




	\begin{enumerate}
		\item 
Les tirs sont indépendants et à chaque tir la probabilité de le réussir est égale à 0,35 : la variable aléatoire $X$ égale au nombre de réussites suit donc une loi binomiale de paramètres $n = 10$ et $p  = 0,35$.
		\item 
		On a $E = np = 10 \times 0,35 = 3,5$ : en moyenne sur $20$ tirs Stéphanie en réussira $7$.
		\item 
		La calculatrice donne $p(X \leqslant 6) \approx 0,97$.
		\item 
		On a de même $p(X \geqslant 6) \approx 0,095$.
	\end{enumerate}
\item 






La probabilité de rater $n$ tirs à 3 points est égale à $0,65^n$, donc elle d'en réussir un est $1 - 0,65^n$.

Il faut donc résoudre dans $\N$, l'inéquation $1 - 0,65^n \geqslant 0,99 \iff 0,01 \geqslant 0,65^n$, soit par croissance de la fonction logarithme népérien:

$\ln 0,01 \geqslant n\ln 0,62 \iff \dfrac{\ln 0,01}{\ln 0,65} \leqslant n$ (car $\ln 0,65 < 0$).

Or $\dfrac{\ln 0,01}{\ln 0,65} \approx 10,69$ ; le plus petit naturel solution est donc 11 : Stéphanie doit tenter 11 tirs.
\end{enumerate}

\bigskip

\end{document}
